\akapit Badane charakterystyki mają przebieg zgodny z teorią: ciąg i moment rosną z kwadratem prędkości obrotowej, a współczynniki $C_T$ i $C_N$ maleją ze wzrostem $\lambda$. Wyniki są zbliżone do danych katalogowych i modelu Quanser, więc pomiary uznajemy za poprawne.

\akapit Obniżenie napięcia zasilania (symulacja rozładowanej baterii) praktycznie nie zmienia zależności $T(n)$ i $M(n)$, ale ogranicza maksymalne obroty, a przez to dostępny ciąg. Napływ powietrza w tunelu zmniejsza ciąg przy tych samych obrotach, ponieważ maleje kąt natarcia łopat.

\akapit Sprawność zespołu napędowego jest najwyższa przy średnich obrotach ($\eta_{zn} \approx 0{,}35$). Wysoka wartość sprawności śmigła dla największego $\lambda$ jest podejrzana, bo moment był tam bardzo mały i niewielki błąd pomiaru mocno zmienia wynik. Błędy mogły wynikać też z dokładności belek tensometrycznych oraz odczytu obrotów z oscyloskopu.
